%==============================================================================
\section{Experiments}
\label{sec:experiments}
%==============================================================================

\subsection{Experimental Setup}
\label{sec:setup}

\paragraph{Dataset.}
We construct a code retrieval dataset from SWE-smith~\cite{yang2025swesmith}, a large-scale collection of synthetic Python bug-fix instances derived from real-world GitHub repositories.
After excluding repositories with failed clones, 44,173 instances across 123 repositories remain.
Our structural extraction pipeline (\S\ref{sec:extraction}) produces inter-procedural data-flow graphs for \placeholder{39,750} instances (\placeholder{90\%}).
Per-seed stratified splits produce $\sim$11,300 training and $\sim$3,000 validation instances, preserving repository-level distribution across partitions.
A separate query-split evaluation set of \placeholder{504} held-out queries is reserved for the D062 protocol (\S\ref{sec:query-split}).

\paragraph{Encoders.}
We evaluate three encoders: Qwen2.5-Coder-3B~\cite{hui2024qwen25coder} (decoder-only, saturated geometry), StarCoder2-3B~\cite{lozhkov2024starcoder2} (decoder-only, healthy separation), and UniXcoder-base~\cite{guo2022unixcoder} (125M encoder-only, architecture-scale boundary).
Full specifications appear in \S\ref{sec:training}.

\paragraph{Training.}
All fine-tuning details follow \S\ref{sec:training}: LoRA ($r{=}16$, $\alpha{=}32$) on $W_q/W_v$, InfoNCE with learnable $\tau$, AdamW ($\text{lr}{=}2{\times}10^{-5}$).
Seed 42 trains for 3 epochs (overfitting diagnostic); seeds 43--44 use 2 epochs with early stopping (patience 2).
We maintain dual checkpoints: $\theta^*_{\text{loss}}$ (best validation loss) and $\theta^*_{\text{R@1}}$ (best validation Recall@1).

\paragraph{Evaluation.}
We report two evaluation protocols:
\begin{itemize}
\item \textbf{Pool-50 (training eval):} Standard pool-based retrieval, pool size 50 (1 positive + 49 negatives from same repository). Reported for comparability with prior work.
\item \textbf{Query-split (D062 protocol):} Held-out queries evaluated at pool sizes \{50, 200, 500\}. Both $\theta^*_{\text{loss}}$ and $\theta^*_{\text{R@1}}$ checkpoints evaluated per configuration. Primary metric: R@1.
\end{itemize}
Bootstrap 95\% confidence intervals computed over query-level outcomes.

\paragraph{Configurations.}
Each encoder $\times$ strategy combination is trained with 3 random seeds (42, 43, 44), yielding:
\begin{itemize}
\item Qwen: 4 strategies $\times$ 3 seeds = 12 runs
\item StarCoder: 4 strategies $\times$ 3 seeds = 12 runs
\item UniXcoder: 4 strategies $\times$ 3 seeds = 12 runs (pool=8 for s42; pool=50 for s43--44)
\end{itemize}
Total: 36 runs. D062 query-split evaluation: \placeholder{14} checkpoint evaluations across strategies and criteria.


\subsection{Main Results}
\label{sec:results}

% ===========================================================
% VERSION-DEPENDENT: Lead table and narrative differ
% ===========================================================

\begin{table*}[t]
\centering
\caption{Pool-50 training evaluation: R@1 (mean $\pm$ std across 3 seeds) by strategy and encoder. $\theta^*_{\text{R@1}}$ checkpoint.}
\label{tab:pool50}
\small
\begin{tabular}{lcccc}
\toprule
Strategy & Qwen2.5-Coder-3B & StarCoder2-3B & UniXcoder-base \\
\midrule
Random           & \placeholder{X.XX $\pm$ X.XX} & \placeholder{X.XX $\pm$ X.XX} & \placeholder{X.XX $\pm$ X.XX} \\
Embedding-dist.  & \placeholder{X.XX $\pm$ X.XX} & \placeholder{X.XX $\pm$ X.XX} & \placeholder{X.XX $\pm$ X.XX} \\
BM25             & \placeholder{X.XX $\pm$ X.XX} & \placeholder{X.XX $\pm$ X.XX} & \placeholder{X.XX $\pm$ X.XX} \\
Dataflow-aware   & \placeholder{X.XX $\pm$ X.XX} & \placeholder{X.XX $\pm$ X.XX} & \placeholder{X.XX $\pm$ X.XX} \\
\bottomrule
\end{tabular}
\end{table*}

\begin{table*}[t]
\centering
\caption{Query-split evaluation (D062 protocol): R@1 by strategy, checkpoint criterion, and pool size. Qwen2.5-Coder-3B, \placeholder{N} seeds.}
\label{tab:querysplit}
\small
\begin{tabular}{llccc}
\toprule
Strategy & Criterion & Pool-50 & Pool-200 & Pool-500 \\
\midrule
Random          & $\theta^*_{\text{loss}}$ & \placeholder{XX.X} & \placeholder{XX.X} & \placeholder{XX.X} \\
Random          & $\theta^*_{\text{R@1}}$  & \placeholder{XX.X} & \placeholder{XX.X} & \placeholder{XX.X} \\
\midrule
Embedding-dist. & $\theta^*_{\text{loss}}$ & \placeholder{XX.X} & \placeholder{XX.X} & \placeholder{XX.X} \\
Embedding-dist. & $\theta^*_{\text{R@1}}$  & \placeholder{XX.X} & \placeholder{XX.X} & \placeholder{XX.X} \\
\midrule
Dataflow-aware  & $\theta^*_{\text{loss}}$ & \placeholder{XX.X} & \placeholder{XX.X} & \placeholder{XX.X} \\
Dataflow-aware  & $\theta^*_{\text{R@1}}$  & \placeholder{XX.X} & \placeholder{XX.X} & \placeholder{XX.X} \\
\bottomrule
\end{tabular}
\end{table*}

\placeholder{[Main results narrative — version-dependent. Version A leads with criterion effect; Version B leads with strategy effect at realistic pool size.]}


\subsection{Analysis}
\label{sec:analysis}

\subsubsection{Checkpoint Selection Criterion}
\label{sec:criterion}

\begin{figure}[t]
\centering
\includegraphics[width=\columnwidth]{\placeholder{figures/fig3_criterion_effect.pdf}}
\caption{Criterion effect vs.\ strategy effect. The R@1 gap between $\theta^*_{\text{R@1}}$ and $\theta^*_{\text{loss}}$ (criterion effect, \placeholder{+16.3pp}) substantially exceeds the gap between strategies under the same criterion (strategy effect, \placeholder{+X.Xpp}).}
\label{fig:criterion-effect}
\end{figure}

Among \placeholder{6/14} completed query-split evaluations, the criterion effect averages \placeholder{+16.3pp}: selecting checkpoints by validation R@1 rather than validation loss yields dramatically higher retrieval performance on held-out queries.
This effect is consistent across strategies, indicating that loss--R@1 divergence (\S\ref{sec:training-dynamics}) is a general property of contrastive fine-tuning rather than a strategy-specific phenomenon.

\placeholder{[Full 14/14 statistics after D062 completion. Report: criterion main effect, strategy main effect, interaction, per-encoder breakdown.]}

\subsubsection{Pool Size Sensitivity}
\label{sec:pool-size}

\begin{figure}[t]
\centering
\includegraphics[width=\columnwidth]{\placeholder{figures/fig5_pool_sensitivity.pdf}}
\caption{R@1 vs.\ pool size by strategy ($\theta^*_{\text{R@1}}$ criterion). Strategy differences \placeholder{amplify / compress} at larger pool sizes.}
\label{fig:pool-sensitivity}
\end{figure}

\placeholder{[Pool size analysis. Report R@1 × strategy × pool size {50, 200, 500}. Key question: do strategy differences amplify or compress at realistic pool sizes?]}

\subsubsection{Training Dynamics}
\label{sec:training-dynamics}

We define two metrics to characterize training trajectories:
\begin{itemize}
\item \textbf{R@1 decay rate} ($\delta_{\text{R@1}}$): per-500-step R@1 drop after peak.
\item \textbf{Loss--R@1 divergence}: whether $\theta^*_{\text{loss}}$ and $\theta^*_{\text{R@1}}$ select different steps.
\end{itemize}

Loss--R@1 divergence occurs in \placeholder{60\%} of completed runs across both strategies, confirming it is a general property of contrastive fine-tuning (\S\ref{sec:criterion}).
Post-peak decay differs by strategy: embedding-distance negatives show \placeholder{$-$2.2\%} per 500 steps vs.\ \placeholder{$-$6.2\%} for random (Qwen seed 44 controlled comparison), suggesting harder negatives provide a regularizing effect.

\placeholder{[Aggregate statistics across all runs after M2 completion.]}

\subsubsection{Negative Difficulty Analysis}
\label{sec:neg-difficulty}

\begin{figure}[t]
\centering
\includegraphics[width=\columnwidth]{figures/fig4_rep_geometry.pdf}
\caption{Representation geometry across encoders and strategies. (a) Positive-negative margin comparison. (b) Similarity distribution. Qwen's near-saturated space (pos\_sim $=$ 0.986) compresses margins to 0.017--0.074; StarCoder maintains 2.6--4.1$\times$ larger margins.}
\label{fig:rep-geometry}
\end{figure}

Table~\ref{tab:neg-difficulty} reports pretrained embedding space statistics.
Both encoders follow the expected difficulty ordering: embedding-distance negatives are hardest (smallest margin), dataflow-aware intermediate, random easiest.
Qwen's margin compression is more severe (77\% reduction from random to embedding) than StarCoder's (65\%), consistent with the representation geometry hypothesis (\S\ref{sec:discuss-when}).

\begin{table}[t]
\centering
\caption{Negative difficulty in pretrained embedding space (2048-token context).}
\label{tab:neg-difficulty}
\small
\begin{tabular}{lcccccc}
\toprule
 & \multicolumn{3}{c}{Qwen2.5-Coder-3B} & \multicolumn{3}{c}{StarCoder2-3B} \\
\cmidrule(lr){2-4} \cmidrule(lr){5-7}
Strategy & pos\_sim & neg\_sim & margin & pos\_sim & neg\_sim & margin \\
\midrule
Random          & 0.986 & 0.911 & 0.074 & 0.799 & 0.605 & 0.194 \\
Embedding-dist. & 0.986 & 0.969 & 0.017 & 0.808 & 0.741 & 0.067 \\
BM25            & \placeholder{X.XX} & \placeholder{X.XX} & \placeholder{X.XX} & \placeholder{X.XX} & \placeholder{X.XX} & \placeholder{X.XX} \\
Dataflow-aware  & 0.986 & 0.956 & 0.029 & 0.801 & 0.682 & 0.119 \\
\bottomrule
\end{tabular}
\end{table}

\subsubsection{Cross-Encoder Comparison}
\label{sec:cross-encoder}

\paragraph{Architecture-scale boundary.}
UniXcoder-base (125M) shows negligible strategy differentiation (\placeholder{0.37pp} span, seed 42) compared to the 3B encoders.
This null result establishes a boundary condition: mining strategy benefits appear contingent on sufficient encoder capacity.

\placeholder{[UniXcoder s43/s44 multi-seed confirmation. Report mean ± std. If null persists → boundary confirmed. If effect emerges → pool-size interaction.]}

\paragraph{Decoder-only encoder comparison.}
\placeholder{[Qwen vs StarCoder strategy rankings. Report whether ordering is consistent across encoders.]}

%==============================================================================
